\documentclass[12pt, a4paper]{article}
\usepackage[utf8]{inputenc}
\usepackage{geometry}
\geometry{a4paper, margin=1in}
\usepackage{tabularx}
\usepackage{booktabs}
\usepackage{hyperref}
\usepackage{enumitem}
\usepackage{titlesec}

% Adjust section spacing
\titlespacing*{\section}{0pt}{1.5ex plus 1ex minus .2ex}{1ex plus .2ex}

\title{\vspace{-2cm}\textbf{Department of Artificial Intelligence and Machine Learning}\\
\large Course Code: AI362IA \hspace{1cm} Course: Big Data Technologies}
\date{}

\begin{document}

\maketitle
\vspace{-1.5cm}

\section*{Project Synopsis}

\begin{table}[h!]
\centering
\renewcommand{\arraystretch}{1.5}
\begin{tabularx}{\textwidth}{|l|X|}
\hline
\textbf{Project Title} & \textbf{Smart Log Analyzer for Scalable System Monitoring} \\ \hline
\textbf{Student Name/s} & Abhilash Maiya Y, Biradar Abhishek Mallikarjun, Harsh Agarwal \\ \hline
\textbf{Student USN/s} & 1RV23AI004, 1RV23AI027, 1RV23AI036 \\ \hline
\end{tabularx}
\end{table}

\section*{Introduction}
Modern digital systems --- including cloud servers, enterprise applications, and distributed platforms --- continuously generate massive volumes of log data. These logs contain critical information about system performance, user activity, security events, and application failures. Traditional log monitoring tools lack the scalability needed to process such high-volume, high-velocity data efficiently. The proposed \textbf{Smart Log Analyzer} leverages Big Data technologies --- specifically \textbf{HDFS}, \textbf{Apache Spark}, and \textbf{Apache Hive} --- to build a distributed, fault-tolerant framework capable of real-time anomaly detection, failure identification, and user behavior analysis across large-scale log datasets.

\section*{Research Gaps}
Existing log monitoring and analysis systems present several unresolved limitations that motivate this work:
\begin{enumerate}[label=\textbf{\arabic*.}]
    \item \textbf{Lack of Scalability in Traditional Tools:} Conventional log analyzers (e.g., grep-based or single-node ELK stacks) fail to handle logs at petabyte scale due to centralized architecture and limited horizontal scalability.
    \item \textbf{Absence of Fault-Tolerant Distributed Storage:} Most existing solutions do not leverage distributed file systems, making them susceptible to data loss and single points of failure in large-scale deployments.
    \item \textbf{Insufficient Real-Time Parallel Processing:} Current tools predominantly rely on sequential batch processing, creating latency bottlenecks when dealing with streaming or continuously generated log data across distributed nodes.
    \item \textbf{Limited Anomaly Detection Capability:} Existing systems largely depend on rule-based threshold alerts, lacking the ability to detect complex, multi-dimensional anomalies and behavioral patterns in log data using distributed computation.
    \item \textbf{Inadequate Structured Query Support for Log Analytics:} There is a gap in integrating SQL-like querying capabilities over raw, unstructured distributed log data, making ad-hoc investigative analysis by system administrators difficult and time-consuming.
\end{enumerate}

\section*{Objectives}
\begin{enumerate}[label=\textbf{\arabic*.}]
    \item To design and implement a distributed log storage framework using \textbf{HDFS} for fault-tolerant, scalable ingestion of raw server and application log files.
    \item To perform parallel in-memory log processing using \textbf{Apache Spark} for efficient filtering, aggregation, and transformation of large-scale log datasets.
    \item To detect system errors, failures, and anomalies by analyzing error-rate trends and irregular access patterns within the distributed log corpus.
    \item To identify high-frequency users and suspicious IP addresses by computing request frequency distributions across the processed log data.
    \item To enable structured analytical querying over processed logs using \textbf{Apache Hive}, facilitating SQL-based reporting on system behavior, traffic patterns, and failure incidents.
\end{enumerate}

\section*{Proposed Methodology}
The system adopts a distributed Big Data processing pipeline built on the Hadoop ecosystem. Raw server and application log files are first ingested and stored in \textbf{HDFS}, ensuring fault-tolerant, block-level distributed storage across cluster nodes. The stored log data is subsequently processed by \textbf{Apache Spark}, which performs parallel in-memory computations to filter error events, detect anomalies, and aggregate access frequencies. The processed results are then loaded into \textbf{Apache Hive}, where structured HiveQL queries are executed to generate analytical reports on system failures, user behavior, and traffic trends. The complete pipeline simulates an enterprise-grade log analytics architecture suitable for operational intelligence and scalable system monitoring.

\section*{Major Tools Used}

\begin{table}[h!]
\centering
\renewcommand{\arraystretch}{1.5}
\begin{tabularx}{\textwidth}{|l|X|}
\hline
\textbf{Tool / Technology} & \textbf{Role in the Project} \\ \hline
\textbf{Apache Hadoop (HDFS)} & Distributed, fault-tolerant storage of raw and processed log files \\ \hline
\textbf{Apache Spark} & Parallel in-memory processing, anomaly detection, and log aggregation \\ \hline
\textbf{Apache Hive} & Structured SQL-like querying and report generation over distributed log datasets \\ \hline
\textbf{Apache Flume} & Reliable ingestion and streaming of log data into HDFS \\ \hline
\textbf{Python / PySpark} & Scripting for data transformation, pattern analysis, and visualization \\ \hline
\textbf{Hadoop MapReduce} & Batch processing for frequency computation and large-scale log summarization \\ \hline
\end{tabularx}
\end{table}

\section*{Expected Outcomes}
\begin{enumerate}[label=\textbf{\arabic*.}]
    \item A fully operational \textbf{distributed log storage system} built on HDFS, capable of storing and managing large-scale log datasets across multiple nodes with fault tolerance.
    \item A \textbf{parallel log processing engine} using Apache Spark that efficiently filters, aggregates, and analyzes log records through in-memory distributed computation.
    \item An \textbf{error and anomaly detection module} that identifies system failures, unusual traffic spikes, and irregular behavioral patterns from log data in near real-time.
    \item A \textbf{user and IP access pattern analyzer} that computes request frequency distributions to flag high-frequency or potentially malicious actors within the log corpus.
    \item A \textbf{Hive-based query and reporting system} that enables administrators to run structured analytical queries and generate actionable insights on system performance and security events.
    \item A \textbf{scalable Big Data monitoring architecture} demonstrating the practical integration of HDFS, Spark, and Hive as an enterprise-grade log analytics pipeline.
\end{enumerate}

\vspace{0.5cm}
\noindent\textit{Department of Artificial Intelligence and Machine Learning --- RV College of Engineering, Bengaluru}

\end{document}
